\section{The matrix energy and its curvature}\label{sec:matrix}

We first record the finite-dimensional facts that will measure the
energy increase under a dyadic split. All spaces and matrices in the
proof are real; \(R^*\) denotes the transpose of \(R\).
For a symmetric positive semidefinite matrix \(P\), define
\[
 \Phi(P)=\tr(P^{3/2}),\qquad \supp P=(\ker P)^\perp.
\]
A symmetric matrix \(K\) is \emph{supported on \(P\)} if it annihilates
\(\ker P\). Symmetry then also puts its range in \(\supp P\).
For such \(K\), define
\begin{equation}\label{eq:hessian}
 b_P(K)=\frac32\sum_{i,j}
          \frac{\abs{K_{ij}}^2}{\sqrt{p_i}+\sqrt{p_j}},
\end{equation}
where \(p_i>0\) are the eigenvalues of \(P\) on its support, and the
entries are taken in an orthonormal eigenbasis of that support.
If \(P=0\), the only supported \(K\) is zero and \(b_0(0)=0\).
Simultaneous orthogonal conjugation and extension by zero preserve
\(\Phi\) and \(b\).

The next two lemmas are the real-matrix versions of the energy
derivative and order comparison results in the companion manuscript
\cite[Section~2]{OpenAI2026}. We give the proofs, including the support
conventions needed for singular Gram matrices.

\begin{lemma}[Hessian of the energy]\label{lem:hessian}
On matrices positive definite on a fixed support, \(\Phi\) is smooth
and, for a symmetric direction \(K\) on that support,
\[
 d\Phi_P(K)=\frac32\tr(\sqrt P\,K),
 \qquad d^2\Phi_P(K,K)=b_P(K).
\]
In particular, \(\Phi\) is convex on the positive semidefinite cone.
\end{lemma}

\begin{proof}
Work on the support and set \(B=\sqrt P\). The derivative of the
squaring map at \(B\) is
\[
 \mathcal L_B(X)=BX+XB.
\]
This is the square-root instance of the classical matrix derivative
calculus \cite[Section~1, Equation~(1)]{Daletskii1957}.
In an eigenbasis of \(P\), this map multiplies entry \((i,j)\) by
\(\sqrt{p_i}+\sqrt{p_j}>0\). The inverse function theorem therefore
gives a smooth square root whose derivative \(X\) in direction \(K\)
satisfies \(BX+XB=K\). Differentiating \(\tr(B^3)\) and using cyclicity,
\[
 d\Phi_P(K)=3\tr(B^2X)=\frac32\tr(BK).
\]
A second differentiation gives
\[
 d^2\Phi_P(K,K)
 =\frac32\inner{K}{\mathcal L_B^{-1}K}_{\HS}
 =b_P(K),
\]
where \(\inner{X}{Y}_{\HS}=\tr(X^*Y)\).
The Hessian is nonnegative. This proves convexity on positive
definite matrices; replacing each matrix by itself plus \(\eta\Id\)
and letting \(\eta\downarrow0\) proves convexity on the closed cone.
\end{proof}

\begin{lemma}[Order comparison]\label{lem:order}
Suppose \(P,Q\) are positive semidefinite matrices on the same space,
\(a\ge1\), and \(P\le aQ\). Every symmetric \(K\) supported on \(P\)
is supported on \(Q\), and
\begin{equation}\label{eq:order}
 b_P(K)\ge a^{-1/2}b_Q(K).
\end{equation}
\end{lemma}

\begin{proof}
For \(x\in\ker Q\), the inequalities
\(0\le\inner{Px}{x}\le a\inner{Qx}{x}=0\) imply \(Px=0\).
Thus \(\ker Q\subseteq\ker P\), which proves the support assertion.

First suppose \(P,Q\) are positive definite. Inversion reverses
positive definite order: conjugating \(0<X\le Y\) by \(X^{-1/2}\),
inverting the resulting matrix at least \(\Id\), and conjugating back
gives \(Y^{-1}\le X^{-1}\). The spectral integral
\[
 \sqrt X=\frac1\pi\int_0^\infty
        X(X+t\Id)^{-1}\,\frac{\mathrm dt}{\sqrt t}
\]
then proves monotonicity of the positive square root
\cite[Theorem, p.~309]{Pedersen1972}, since
\(X(X+t\Id)^{-1}=\Id-t(X+t\Id)^{-1}\). Consequently
\(\sqrt P\le\sqrt a\,\sqrt Q\).

For a positive semidefinite matrix \(A\), left and right multiplication
by \(A\) are positive semidefinite self-adjoint operators on the full
Hilbert--Schmidt matrix space: their quadratic forms are
\(\tr(X^*AX)\) and \(\tr(X^*XA)\). Hence
\[
 \mathcal L_{\sqrt P}\le\sqrt a\,\mathcal L_{\sqrt Q},
 \qquad
 \mathcal L_{\sqrt P}^{-1}
       \ge a^{-1/2}\mathcal L_{\sqrt Q}^{-1}.
\]
Taking the quadratic form at \(K\) and using Lemma~\ref{lem:hessian}
proves \eqref{eq:order}.

In the semidefinite case, \(a\ge1\) gives
\(P+\eta\Id\le a(Q+\eta\Id)\). Apply the positive definite result
and let \(\eta\downarrow0\). In a full eigenbasis of either matrix,
every entry of \(K\) meeting its kernel is zero. Formula
\eqref{eq:hessian} therefore gives convergence to the two supported
quadratic forms in \eqref{eq:order}.
\end{proof}

\begin{corollary}[A midpoint gap]\label{cor:gap}
If \(T\ge0\) and \(U=U^*\) satisfy \(-T\le U\le T\), then \(U\) is
supported on \(T\) and
\begin{equation}\label{eq:gap}
 \frac{\Phi(T+U)+\Phi(T-U)}2-\Phi(T)
       \ge\frac{1}{2\sqrt2}b_T(U).
\end{equation}
\end{corollary}

\begin{proof}
For \(x\in\ker T\), positivity of \(T+U\) and \(T-U\), whose quadratic
forms both vanish at \(x\), implies \(Ux=0\).
For \(\abs t<1\),
\[
 (1-\abs t)T\le T+tU\le(1+\abs t)T\le2T.
\]
Thus \(T+tU\) has the same support as \(T\). The function
\(f(t)=\Phi(T+tU)\) satisfies, by Lemmas~\ref{lem:hessian}
and~\ref{lem:order},
\[
 f''(t)=b_{T+tU}(U)\ge 2^{-1/2}b_T(U).
\]
Subtracting \(b_T(U)t^2/(2\sqrt2)\) leaves a convex function on
\((-1,1)\), continuous on \([-1,1]\). Its midpoint inequality is
\eqref{eq:gap}. If \(T=0\), all terms are zero.
\end{proof}
